\chapter{Энергия Гиббса и равновесие}\label{ch:gibbs}

\section{Что стремится к минимуму}

Изолированная система движется к максимуму энтропии: это второе начало термодинамики. Но кусок металла в печи
не изолирован --- он обменивается теплом с окружением при постоянной температуре $T$ и постоянном давлении $p$.
Если посчитать энтропию \emph{системы вместе с печью}, окажется, что её рост равносилен убыванию одной величины,
относящейся только к образцу, --- \textbf{энергии Гиббса}
\begin{equation}
G = H - TS,\qquad H = U + pV .
\end{equation}
Здесь $H$ --- энтальпия (для твёрдых тел почти то же, что внутренняя энергия: член $pV$ при атмосферном давлении
мал), $S$ --- энтропия.

\begin{idea}
Почему именно $H - TS$? Когда образец выделяет теплоту $Q$, энтропия печи растёт на $Q/T$. Выделение теплоты
означает уменьшение $H$ образца. Значит, суммарная энтропия растёт, если $\Delta S_{\text{обр}} - \Delta H/T > 0$,
то есть если $\Delta(H - TS) < 0$. Энергия Гиббса --- это энтропия Вселенной, пересчитанная на образец и взятая
со знаком минус.
\end{idea}

Два слагаемых тянут в разные стороны. Энтальпия любит порядок: атомы ставятся так, чтобы связи были
прочнее. Энтропия любит беспорядок: много способов расставить атомы. Множитель $T$ решает, кто сильнее. При
низкой температуре выигрывает энтальпия --- плотные упорядоченные фазы, соединения, выпадение второй фазы. При
высокой --- энтропия: растворы, расплавы, высокотемпературные модификации.

\section{Устойчивость, метастабильность и барьер}

Фаза устойчива при данных $T$ и $p$, если у неё наименьшая $G$ из всех доступных состояний. Однако
\emph{наименьшая} $G$ --- не единственное, что важно. На рис.~\ref{fig:f1_gibbs}б показан профиль энергии по
некоторой «координате» перестройки атомов. Состояние в локальной яме \textbf{метастабильно}: малые отклонения
возвращают его назад, и, чтобы попасть в более глубокую яму, система должна преодолеть \textbf{барьер}
$\Delta G_a$. Состояние на вершине \textbf{неустойчиво} --- любое отклонение уводит его вниз.

\fig{f1_gibbs}{а) Энергии Гиббса двух фаз в зависимости от температуры. Кривые пересекаются при $T_e$; ниже
устойчива $\alpha$, выше --- $\beta$. На расстоянии $\Delta T$ от $T_e$ разность $\Delta G$ --- движущая сила
превращения. б) Устойчивое, метастабильное и неустойчивое состояния; барьер $\Delta G_a$ держит метастабильное
состояние.}

Почти всё интересное в материаловедении --- метастабильно. Пересыщенный твёрдый раствор водорода в цирконии
после быстрого охлаждения --- метастабилен: гидрид выгоднее, но чтобы он появился, нужно преодолеть барьер
зарождения (глава~\ref{ch:nucleation}). Скорость, с которой система переходит барьер, пропорциональна $\exp(-\Delta G_a/k_BT)$.
Отсюда два вывода, которые будут повторяться во всём тексте:
\begin{enumerate}
\item термодинамика говорит, \emph{куда} может идти процесс (знак $\Delta G$);
\item кинетика --- \emph{как быстро} (высота барьера и подвижность атомов).
\end{enumerate}

\section{Движущая сила вблизи температуры превращения}

Пусть две фазы равновесны при $T_e$: $G_\alpha(T_e) = G_\beta(T_e)$. Немного ниже, при $T = T_e - \Delta T$,
разность энергий
\[
\Delta G = \Delta H - T\,\Delta S .
\]
При $T_e$ она равна нулю, значит $\Delta S = \Delta H/T_e$. Если $\Delta H$ и $\Delta S$ слабо зависят от
температуры (так и есть в узком интервале), то
\begin{keyf}
\[
\Delta G \approx \Delta H - T\,\frac{\Delta H}{T_e} = \frac{\Delta H\,\Delta T}{T_e}.
\]
\end{keyf}
Движущая сила растёт пропорционально переохлаждению. На самой температуре равновесия её нет вовсе --- и поэтому
никакое превращение не идёт \emph{точно} при $T_e$: нужно переохлаждение (или перегрев), чтобы появилась сила,
способная оплатить барьер. Это первая встреча с идеей, которая объяснит гистерезис растворения и выпадения
гидридов (глава~\ref{ch:solutions}).

\section{Химический потенциал}

В сплаве атомы разных сортов могут переходить из фазы в фазу. Нужна величина, которая говорит, выгодно ли
атому сорта $i$ перейти. Это \textbf{химический потенциал}
\begin{equation}
\mu_i = \left(\frac{\partial G}{\partial n_i}\right)_{T,p,n_{j\ne i}},
\end{equation}
изменение энергии Гиббса при добавлении одного моля компонента $i$ (при неизменном прочем).

\begin{idea}
Химический потенциал --- это «давление» для атомов данного сорта. Атомы сами собой перетекают туда, где
их $\mu$ меньше, как газ --- туда, где давление ниже. Равновесие наступает, когда $\mu_i$ одинаков во всех фазах
и во всех точках. Из этого простого условия следуют и общая касательная (глава~\ref{ch:solutions}), и диффузия к растянутым местам
(глава~\ref{ch:diffusion}), и сдвиг растворимости под нагрузкой (глава~\ref{ch:nucleation}).
\end{idea}

Условие равновесия двух фаз $\alpha$ и $\beta$ по всем компонентам:
\begin{equation}
\mu_i^\alpha = \mu_i^\beta\qquad\text{для каждого } i .
\end{equation}
В двухкомпонентной системе это два уравнения. Графически --- общая касательная к кривым $G(X)$ двух фаз; к ней мы
придём в следующей главе.

\section{Масштабы энергий}

Полезно заранее держать в голове порядки величин (на моль):
\begin{center}\small
\begin{tabular}{lr}
\toprule
энергия связи атомов металла (теплота сублимации Zr) & $\sim 600$~кДж/моль \\
теплота растворения водорода в $\alpha$-Zr (из наклона TSSD, глава~\ref{ch:solutions}) & $\sim 40$~кДж/моль \\
тепловая энергия $RT$ при 300\degC & 4.8~кДж/моль \\
движущая сила выпадения гидрида на линии выпадения & $\sim 5$~кДж/моль H \\
упругая энергия тонкой пластинки гидрида (глава~\ref{ch:elastic}) & $\sim 3$~кДж/моль H \\
работа нагрузки 100~МПа на несоответствии гидрида & $\sim 0.07$~кДж/моль H \\
\bottomrule
\end{tabular}
\end{center}
Последние три строки --- суть всей задачи о переориентации: работа нагрузки в десятки раз меньше движущей силы и
упругой энергии, и всё же она решает, как ляжет пластинка. Почему так получается --- главы~\ref{ch:elastic} и 6.

\begin{exercise}
Две фазы равновесны при $T_e = 1000$~К, теплота превращения $\Delta H = 10$~кДж/моль. Оцените движущую силу при
переохлаждении на 10 и на 100~К. Во сколько раз она меньше $RT$ в каждом случае?
\end{exercise}

\begin{exercise}
Скорость перехода через барьер пропорциональна $e^{-\Delta G_a/RT}$. Барьер 120~кДж/моль. Во сколько раз быстрее
процесс при 400\degC, чем при 300\degC?
\end{exercise}

\begin{exercise}
Покажите, что из условия $\dd G = V\dd p - S\dd T$ для двух равновесных фаз следует уравнение Клапейрона
$\dd p/\dd T = \Delta H/(T\Delta V)$. Как, по аналогии, напряжение $\sigma$, совершающее работу на деформации
превращения $\eps^*$, сдвигает температуру равновесия? (К этому вопросу мы вернёмся в главе~\ref{ch:nucleation}.)
\end{exercise}
